%\documentclass[12pt,letterpaper]{article}



%%%
%% Required packages:
%%
\usepackage{etex}
\usepackage{amsmath}
\usepackage{amssymb}
\usepackage{amstext}
\usepackage{amsthm}
\usepackage{fancyhdr,fancybox}
\usepackage{mathrsfs} % need for \mathscr command
\usepackage{multicol}
\usepackage[normalem]{ulem}
%\usepackage{pictex,pictexplus}
% \usepackage{pictexwd}
\usepackage{rotating}
\usepackage{url}
\usepackage{xfrac}
\usepackage{booktabs}
\usepackage{color,soul}
\usepackage{comment}

\usepackage{hyperref}
\hypersetup{
    colorlinks=true,     % false: boxed links; true: colored links
    linkcolor=blue,      % color of internal links
    citecolor=blue,      % color of links to bibliography
    filecolor=blue,      % color of file links
    urlcolor=blue        % color of external links
}


\usepackage{tikz}
\usetikzlibrary{backgrounds,calc}

%%
%% Common formatting commands
%%
\setlength{\parindent}{0in}
\setlength{\textwidth}{7in}
\setlength{\evensidemargin}{-0.25in}
\setlength{\oddsidemargin}{-0.25in}
\setlength{\parskip}{.5\baselineskip}
\setlength{\topmargin}{-0.5in}
\setlength{\textheight}{9in}

%               Problem and Part
\newcounter{problemnumber}
\newcounter{partnumber}
\newcounter{subpartnumber}
\newcommand{\Problem}{%
  \refstepcounter{problemnumber}%
  \setcounter{partnumber}{0}%
  \setcounter{subpartnumber}{0}%
  \item[\makebox{\hfil\textbf{\theproblemnumber.}\hfil}]%
  }
\newcommand{\InProblem}[2][1.6]{%
  \refstepcounter{problemnumber}%
  \setcounter{partnumber}{0}%
  \setcounter{subpartnumber}{0}%
  \textbf{\theproblemnumber.}\ \ \parbox[t]{#1in}{#2}%
}
\newcommand{\InSmallProblem}[1]{\InProblem[1.05]{#1}}
\newcommand{\NoProblem}[1]{%
  \mbox{\hfil\textbf{#1}\hfil}%
  }
\newcommand{\UnProblem}{%
  \refstepcounter{problemnumber}%
  \setcounter{partnumber}{0}%
  \setcounter{subpartnumber}{0}%
  \mbox{\hfil\textbf{\theproblemnumber.}\hfil}%
  }
\newcommand{\InFlexProblem}[1]{%
  \refstepcounter{problemnumber}%
  \setcounter{partnumber}{0}%
  \setcounter{subpartnumber}{0}%
  \textbf{\theproblemnumber.}\ \ #1%
}
%%  Part:
\newcommand{\Part}{%
  \renewcommand{\thepartnumber}{(\alph{partnumber})}%
  \refstepcounter{partnumber}
  \setcounter{subpartnumber}{0}%
  \item[(\alph{partnumber})]%
}
\newcommand{\InPart}[2][1.75]{%
  \renewcommand{\thepartnumber}{(\alph{partnumber})}%
  \refstepcounter{partnumber}%
  \setcounter{subpartnumber}{0}%
  (\alph{partnumber})\ \ \parbox[t]{#1in}{#2}%
}
\newcommand{\UnPart}{%
  \refstepcounter{partnumber}%
  \renewcommand{\thepartnumber}{(\alph{partnumber})}%
  \setcounter{subpartnumber}{0}%
  (\alph{partnumber})%
}
\newcommand{\InLargePart}[1]{\InPart[3.05]{#1}}
\newcommand{\InFullPart}[1]{\InPart[6.2]{#1}}
\newcommand{\InSmallPart}[1]{\InPart[1.05]{#1}}
%% SubPart:
\newcommand{\SubPart}{%
  \refstepcounter{subpartnumber}%
  \item[(\roman{subpartnumber})]%
  }
\newcommand{\InSubPart}[2][2.25]{%
  \stepcounter{subpartnumber}%
  (\roman{subpartnumber})\ \ \parbox[t]{#1in}{#2}%
}
\newcommand{\InSmallSubPart}[1]{\InSubPart[1.5]{#1}}
\newcommand{\InTinySubPart}[1]{%
  \stepcounter{subpartnumber}%
  (\roman{subpartnumber})\ \ \makebox{#1}%
}
\newcommand{\InMySubPart}[2][2.25]{%
  \stepcounter{subpartnumber}%
  \parbox[t]{#1in}{\makebox[0.3in][l]{(\roman{subpartnumber})}\ \ {#2}}%
}
\newcommand{\TrueFalse}{\textbf{T}\ \ \textbf{F}\ \ }
\newcommand{\TFPart}[1]{% 
  \stepcounter{partnumber}%
  \setcounter{subpartnumber}{0}%
  \item[(\alph{partnumber})] \TrueFalse\parbox[t]{5.5in}{#1}}

\newcommand{\Points}[1]{(#1 points) \ }
\newcommand{\Point}[1]{(#1 point) \ }


%% For Answers:
\newcounter{answernumber}
\newcommand{\Answer}{%
  \refstepcounter{answernumber}%
  \setcounter{partnumber}{0}%
  \setcounter{subpartnumber}{0}%
  \item[\makebox{\hfil\textbf{\theanswernumber.}\hfil}]%
  }


% Setting up the headers for the first page
%\chead{} % will default to what is typed in the file.
\fancyhead[L]{\textsc{Math 22a}}
\fancyhead[R]{\textsc{Fall 2024}}
\fancyfoot[R]{
	\vspace*{\fill}\footnotesize\textcopyright{} \emph{Harvard Math 22a}	
}
\renewcommand{\headrulewidth}{0pt}

%\fancyhead[C]{}
%	\chead{}	
%	\lhead{}
%	\rhead{}
%	\cfoot{}


% Putting copyright on the footer of each page
%\fancypagestyle{secondpage}{
%	\fancyhead[C]{TESTing again} % change this to be blank
%		%\chead{}	
%	\fancyhead[L]{bears} % change this to be blank
%	\fancyhead[R]{dogs} % change this to be blank
%	\fancyfoot[R]{ %keeps the footer the same
%		\vspace*{\fill}\footnotesize\textcopyright{} \emph{Harvard Math 22a}	
%	}
%}
%\renewcommand{\headrulewidth}{0pt}
%\pagestyle{fancy}
\pagestyle{empty}


%\lhead{}
%\rhead{}
%\rfoot{\vspace*{\fill}\footnotesize\textcopyright{} \emph{Harvard Math 22a}}
%\renewcommand{\headrulewidth}{0pt}
%\pagestyle{fancy}




%%
%% Common shortcut commands:
%%
\newcommand{\ds}{\displaystyle}
\newcommand{\sgn}{\operatorname{sgn}}
\renewcommand{\Pr}{\operatorname{P}}
\newcommand{\CPr}[3][{}]{\Pr_{\text{#1}}\left(#2\,|\,#3\right)}

\newcommand{\C}{\mathbb{C}}
\newcommand{\R}{\mathbb{R}}
\newcommand{\Z}{\mathbb{Z}}
\newcommand{\N}{\mathbb{N}}
\newcommand{\Q}{\mathbb{Q}}
\newcommand{\D}{\mathbb{D}}

\newcommand{\rref}{\operatorname{rref}}
\newcommand{\im}{\operatorname{im}}
\newcommand{\rank}{\operatorname{rank}}
\newcommand{\diag}{\operatorname{diag}}
\newcommand{\Span}{\operatorname{span}}
%\renewcommand{\vec}[1]{\mathbf{#1}}
\newcommand{\charpoly}{\mathscr{P}}
\newcommand{\nicefrac}[2]{\sfrac{#1}{#2}} % using xfrac package
\newcommand{\topology}{\mathcal{T}}
\newcommand{\Inn}{\operatorname{Inn}}

\newcommand{\blankbox}[2]{\fbox{\rule{#1}{0in}\rule{0in}{#2}}}

\newenvironment{myindentpar}[1]%
 {\begin{list}{}%
         {\setlength{\leftmargin}{#1}
          \setlength{\rightmargin}{\leftmargin}}%
         \item[]%
 }
 {\end{list}}
\newtheorem{theorem}{Theorem}
\newtheorem*{NStheorem}{Theorem}
\newtheorem{cor}[theorem]{Corollary}
\newtheorem*{claim}{Claim}
\newtheorem{defn}{Definition}

\newenvironment{amatrix}[1]{%
  \left(\begin{array}{@{}*{#1}{c}|c@{}}
}{%
  \end{array}\right)
}

%--------grstep
% For denoting a Gauss' reduction step.
% Use as: \grstep{\rho_1+\rho_3} or \grstep[2\rho_5 \\ 3\rho_6]{\rho_1+\rho_3}
\newcommand{\grstep}[2][\relax]{%
   \ensuremath{\mathrel{
       {\mathop{\longrightarrow}\limits^{#2\mathstrut}_{
                                     \begin{subarray}{l} #1 \end{subarray}}}}}}
\newcommand{\swap}{\leftrightarrow}


%\usetikzlibrary{arrows}
%\usepackage{wrapfig}

%\makeatletter
%\renewcommand*\env@matrix[1][*\c@MaxMatrixCols c]{%
%	\hskip -\arraycolsep
%	\let\@ifnextchar\new@ifnextchar
%	\array{#1}}
%\makeatother
%
%\def\env@matrix{\hskip -\arraycolsep
%	\let\@ifnextchar\new@ifnextchar
%	\array{*\c@MaxMatrixCols c}}
%
\section{{Diagonalization Continued, $\vec\S$5.3}}% 
\chead{\Large\textbf{Diagonalization Continued, $\vec\S$5.3}}% 

%\begin{document}
\thispagestyle{fancy}

Recall from last class:

\begin{itemize}
\item If $A$ is an $n \times n$ matrix, then $\lambda \in \mathbb{R}$ is an {\bf eigenvalue} of $A$ if there exists a non-zero vector $v$, called an {\bf eigenvector} corresponding to $\lambda$, such that $$Av=\lambda v.$$
\item If $\lambda$ is an eigenvalue of an $n \times n$ matrix $A$, then the set of all solutions $v$ to $Av = \lambda v$ is the {\bf eigenspace} of $A$ corresponding to $\lambda$.
\item {\bf Theorem 1:} The eigenvalues of a triangular matrix are the diagonal entries.
\item {\bf Theorem 2:} If $v_1, \dots, v_p$ are eigenvectors corresponding to distinct eigenvalues $\lambda_1, \dots, \lambda_p$ of an $n\times n$ matrix $A$, then $v_1, \dots, v_p$ are linearly independent.
\item If $A$ is an $n\times n$ matrix, then the characteristic polynomial of $A$, denoted by $p_A(\lambda)$, is given by $$p_A(\lambda)=\text{det}(A-\lambda I_n).$$
\item $\lambda$ is an eigenvalue of an $n\times n$ matrix $A$ if and only if $\lambda$ is a zero of the characteristic polynomial.
\item If $A$ and $B$ are similar matrices, then they have the same characteristic polynomials and hence the same eigenvalues (with the same algebraic multiplicities). 
\item An $n\times n$ matrix $A$ is {\bf diagonalizable} if $A$ is similar to a diagonal matrix.
\item A basis of $\mathbb{R}^n$ diagonalizing $A$ is called an {\bf eigenbasis}. 
\item {\bf Diagonalization Theorem:} An $n \times n$ matrix $A$ is diagonalizable if and only if $A$ has $n$ linearly independent eigenvectors. Moreover, $A=PDP^{-1}$ with $D$ diagonal if and only if the columns of $P$ are $n$ linearly independent eigenvectors of $A$. In this case the diagonal entires of $D$ are the corresponding eigenvalues of $A$.
\end{itemize}



\newpage



\begin{enumerate}

\Problem Let $A=\begin{bmatrix} 
2 & 4 & 3\\
-4 & -6 & -3\\
3 & 3 & 1\\
\end{bmatrix}$ If possible, diagonalize $A$.


\vfill
\vfill

{\bf Theorem 6:} An $n \times n$ matrix $A$ with $n$ distinct eigenvalues is diagonalizable.

\Problem Prove the theorem.

\vfill

\newpage

\Problem In each part, determine if there is a basis in which the linear transformation is represented by a diagonal matrix.
\Part Let $T: \mathbb{P}_2 \to \mathbb{P}_2$ given by $T(p)= p(t) +(t+1)p'(t)$.

\vfill

\Part Let $T: \mathbb{P}_2 \to \mathbb{P}_2$ given by $T(p) = p(1)+p'(0) t + (p'(0)+p''(0))t^2.$

\vfill




\end{enumerate}

\begin{comment}
\newpage
\chead{\Large\textbf{Diagonalization Again-- Answers / Solutions}}%
\lhead{}
\rhead{}
\thispagestyle{fancy}

\begin{enumerate}


\Answer In order to find the eigenvalues, we compute the zeros of the characteristic polynomial. Note $p_A(\lambda)= \det \begin{bmatrix} 2-\lambda & 4 & 3\\ -4 & -6 - \lambda & -3\\ 3 & 3 & 1-\lambda \end{bmatrix}=-(\lambda-1)(\lambda-2)^2$. Thus the eigenvalues of $A$ are 1 and -2 (with algebraic multiplicity 2). We find the eigenspaces for each now.
  
\underline{$\lambda=1$}. In this case, we get $$E_1 = N(A-I_3) = \text{Span}\left \{\begin{bmatrix} 1\\ -1\\ 1\\  \end{bmatrix} \right\}.$$

\underline{$\lambda=-2$}. In this case, we get $$E_{-2} = N(A+2I_3) = \text{Span}\left \{\begin{bmatrix} -1\\1\\0\\ \end{bmatrix} \right\}.$$

Hence we have only two linearly independent eigenvectors and by the Diagonalization theorem, $A$ cannot be diagonalized.


\Answer 

\begin{proof}
Suppose $A$ is an $n\times n$ matrix with $n$ distinct eigenvalues. Let $v_1, \dots, v_n$ be $n$ eigenvectors each corresponding to a different eigenvalue. We proved last time that eigenvectors corresponding to different eigenvalues must be linearly independent. Thus by the Diagonalization Theorem, $A$ must be diagonalizable. 
\end{proof}
            \Answer In each part we use the standard basis $\mathcal{B}=\{1,t, t^2\}.$
            \begin{enumerate}
            \item Plugging in each basis vector to $T$, we get $$T(1)=1, T(t)= 2t+1, T(t^2)= 3t^2+2t.$$ Expressing each in terms of basis $\mathcal{B}$, we get $[T(1)]_\mathcal{B}=\begin{bmatrix} 1\\ 0\\ 0\\ \end{bmatrix}$, $[T(t)]_\mathcal{B}=\begin{bmatrix} 1\\ 2\\ 0\\ \end{bmatrix}$, and $[T(t^2)]_\mathcal{B}=\begin{bmatrix} 0\\ 2\\ 3\\ \end{bmatrix}$. 
            
            Thus $$[T]_\mathcal{B}=[ [T(1)]_\mathcal{B} \; [T(t)]_\mathcal{B} \; [T(t^2)]_\mathcal{B} ]=\begin{bmatrix}
            1 & 1 & 0\\ 0 & 2 & 2\\ 0 & 0 & 3\\ \end{bmatrix}.$$
            
            Since $[T]_\mathcal{B}$ is upper-triangular, the eigenvalues are the diagonal entries. Thus the eigenvalues of $T$ are 1, 2, and 3. We compute each eigenspace. 
            
            \underline{$\lambda=1$} Since $[T]_\mathcal{B}-I_3=\begin{bmatrix} 0 & 1 & 0\\ 0 & 1 & 2\\ 0 & 0 & 2\\ \end{bmatrix}$, we get $$N([T]_\mathcal{B}-I_3)=\text{Span}\left\{\begin{bmatrix}1\\ 0\\ 0\\ \end{bmatrix} \right\}.$$ 
            
            Thus the eigenspace for $\lambda=1$ is given by $E_1=\text{Span}\{1\}.$
            
            \underline{$\lambda=2$} Since $[T]_\mathcal{B}-2I_3=\begin{bmatrix} -1 & 1 & 0\\ 0 & 0 & 2\\ 0 & 0 & 1\\ \end{bmatrix}$, we get $$N([T]_\mathcal{B}-2I_3)=\text{Span}\left\{\begin{bmatrix}1\\ 1\\ 0\\ \end{bmatrix} \right\}.$$ 

            Thus the eigenspace for $\lambda=2$ is given by $E_2=\text{Span}\{1+t\}.$
            
            \underline{$\lambda=3$} Since $[T]_\mathcal{B}-3I_3=\begin{bmatrix} -2 & 1 & 0\\ 0 & -1 & 2\\ 0 & 0 & 0\\ \end{bmatrix}$, we get $$N([T]_\mathcal{B}-3I_3)=\text{Span}\left\{\begin{bmatrix}1\\ 2\\ 1\\ \end{bmatrix} \right\}.$$  
            
            Thus the eigenspace for $\lambda=3$ is given by $E_3=\text{Span}\{1+2t+t^2\}.$   
            
            Finally we take the basis $\mathcal{C}=\{1,1+t,1+2t+t^2\}$, and then $[T]_\mathcal{C}=\begin{bmatrix} 1 & 0 & 0\\ 0 & 2 &0 \\ 0 & 0 & 3\\ \end{bmatrix}.$  
            
                  
            \item Plugging in each basis vector to $T$, we get $$T(1)=1, T(t)= 1+t+t^2, T(t^2)= 1+2t^2.$$ Expressing each in terms of basis $\mathcal{B}$, we get $[T(1)]_\mathcal{B}=\begin{bmatrix} 1\\ 0\\ 0\\ \end{bmatrix}$, $[T(t)]_\mathcal{B}=\begin{bmatrix} 1\\ 1\\ 1\\ \end{bmatrix}$, and $[T(t^2)]_\mathcal{B}=\begin{bmatrix} 1\\ 0\\ 2\\ \end{bmatrix}$. 
            
            Thus $$[T]_\mathcal{B}=[ [T(1)]_\mathcal{B} \; [T(t)]_\mathcal{B} \; [T(t^2)]_\mathcal{B} ]=\begin{bmatrix}
            1 & 1 & 1\\ 0 & 1 & 0\\ 0 & 1 & 2\\ \end{bmatrix}.$$
            
            The characteristic polynomial of $A$ is $p_A(\lambda)=\det{(A-\lambda I_n)}=-(\lambda-1)^2(\lambda-2).$. The eigenvalues of $A$ are 1 (with algebraic multiplicity 2) and 2 (with algebraic multiplicity 1).
            
            \underline{$\lambda=1$} Since $[T]_\mathcal{B}-I_3=\begin{bmatrix} 0 & 1 & 1\\ 0 & 0 & 0\\ 0 & 1 & 1\\ \end{bmatrix}$, we get $$N([T]_\mathcal{B}-I_3)=\text{Span}\left\{\begin{bmatrix}1\\ 0\\ 0\\ \end{bmatrix},\begin{bmatrix}0\\ 1\\ -1\\ \end{bmatrix} \right\}.$$ 
            
            Thus the eigenspace for $\lambda=1$ is given by $E_1=\text{Span}\{1, t-t^2\}.$
            
            \underline{$\lambda=2$} Since $[T]_\mathcal{B}-2I_3=\begin{bmatrix} -1 & 1 & 1\\ 0 & -2 & 0\\ 0 & 1 & 0\\ \end{bmatrix}$, we get $$N([T]_\mathcal{B}-2I_3)=\text{Span}\left\{\begin{bmatrix}1\\ 0\\ 1\\ \end{bmatrix} \right\}.$$ 
            
            Thus the eigenspace for $\lambda=2$ is given by $E_2=\text{Span}\{1+t^2\}.$
           
            
            Finally we take the basis $\mathcal{C}=\{1,t-t^2,1+t^2\}$, and then $[T]_\mathcal{C}=\begin{bmatrix} 1 & 0 & 0\\ 0 & 1 &0 \\ 0 & 0 & 2\\ \end{bmatrix}.$   
            \end{enumerate}      

\end{enumerate}  

\end{comment}

%\end{document}


%%% Local ables: 
%%% mode: latex
%%% TeX-master: t
%%% End: 
